% SPDX-FileCopyrightText: 2026 AntheaLiles
% SPDX-License-Identifier: CC-BY-SA-4.0
% ==============================================================================
% preamble-article.tex — Préambule LuaLaTeX PDF/UA-2 haute performance pour Org-mode
% Requis dans le manuscrit Org : défini via org-latex-classes dans Emacs
% ==============================================================================

%%% 01. Outils de base & Langue
\usepackage{etoolbox}
\usepackage[main=french,english]{babel}
\usepackage[dvipsnames,svgnames,table]{xcolor}
\definecolor{customgray}{HTML}{505050}

%%% 02. Géométrie & Zone d'annotation
%
% A4 (21 cm) = Marges (2 x 2.4 cm) + Colonne lecture (12.1 cm / ~60 car.) + Marge externe (3.6 cm).
%
\newlength{\zoneannotation}        \setlength{\zoneannotation}{3.6cm}
\newlength{\ecartannotation}       \setlength{\ecartannotation}{0.5cm}
\newlength{\largeurimpression}
\newlength{\debordementannotation}
% Les deux largeurs dérivent de \textwidth, qu'elles doivent lire APRÈS
% geometry : évaluées avant, elles prendraient la largeur par défaut de la
% classe (360 pt en 11 pt, contre 12,1 cm = 344 pt ici) et toute table ou
% figure pleine largeur dépasserait la zone d'annotation d'environ 5,5 mm.
% Elles sont calculées après le chargement de geometry, puis de nouveau à
% \begin{document}, pour suivre un éventuel \geometry{} du document.
\newcommand{\calculerlargeurs}{%
  \setlength{\largeurimpression}{\dimexpr\textwidth+\ecartannotation+\zoneannotation\relax}%
  \setlength{\debordementannotation}{\dimexpr\ecartannotation+\zoneannotation\relax}}

\usepackage[a4paper, twoside,
            top=3.2cm, bottom=3.2cm, inner=2.4cm, outer=2.4cm,
            marginparwidth=\zoneannotation, marginparsep=\ecartannotation,
            includemp]{geometry}
\calculerlargeurs
\AtBeginDocument{\calculerlargeurs}

% Page de garde symétrique sans zone d'annotation
\AtBeginDocument{%
  \let\titreoriginal\maketitle
  \renewcommand{\maketitle}{%
    \newgeometry{a4paper, top=3.2cm, bottom=3.2cm, left=2.4cm, right=2.4cm}%
    \titreoriginal
    \restoregeometry}}

\usepackage{fancyhdr}
\pagestyle{fancy}\fancyhf{}
\renewcommand{\headrulewidth}{0pt}
\setcounter{secnumdepth}{3}
\setlength{\columnsep}{0.8cm}
\setlength{\parskip}{0.5em}
\setlength{\parindent}{0pt}
\raggedbottom

%%% 03. Polices & Typographie (Optimisé LuaLaTeX & Tagged PDF)
%
% Note Perf : protrusion=true est très rapide. expansion=false évite le goulot d'étranglement
% de calcul microtypographique sous LuaTeX combiné au balisage PDF/UA.
%
\usepackage{fontspec}
\usepackage{unicode-math}
\usepackage[protrusion=true,expansion=false,final]{microtype}

\setmainfont{Luciole-Regular.ttf}[
  ItalicFont=Luciole-Regular-Italic.ttf,
  BoldFont=Luciole-Bold.ttf,
  BoldItalicFont=Luciole-Bold-Italic.ttf]
\setmonofont{IosevkaSS05-Regular.ttf}[
  ItalicFont=IosevkaSS05-Italic.ttf,
  BoldFont=IosevkaSS05-Bold.ttf,
  BoldItalicFont=IosevkaSS05-BoldItalic.ttf]
\setmathfont{Luciole-Math.otf}

\tolerance=1500
\hyphenpenalty=100
\emergencystretch=1.5em
\clubpenalty=10000
\widowpenalty=10000
\displaywidowpenalty=10000

\makeatletter
% Hiérarchie de corps
\renewcommand{\tiny}        {\@setfontsize\tiny{7}{8.5}}
\renewcommand{\scriptsize}  {\@setfontsize\scriptsize{8}{10}}
\renewcommand{\footnotesize}{\@setfontsize\footnotesize{9}{11.5}}
\renewcommand{\small}       {\@setfontsize\small{10}{13}}
\renewcommand{\normalsize}  {\@setfontsize\normalsize{11}{14.3}}
\renewcommand{\large}       {\@setfontsize\large{13}{16}}
\renewcommand{\Large}       {\@setfontsize\Large{15}{18}}
\renewcommand{\LARGE}       {\@setfontsize\LARGE{17}{20}}
\renewcommand{\huge}        {\@setfontsize\huge{20}{23}}
\renewcommand{\Huge}        {\@setfontsize\Huge{23}{26}}

% Titres sans gras (préservation des sockets de balisage PDF/UA-2)
\renewcommand{\section}{\@startsection{section}{1}{\z@}{-3.5ex \@plus -1ex \@minus -.2ex}{2.3ex \@plus.2ex}{\normalfont\Large}}
\renewcommand{\subsection}{\@startsection{subsection}{2}{\z@}{-3.25ex\@plus -1ex \@minus -.2ex}{1.5ex \@plus .2ex}{\normalfont\large}}
\renewcommand{\subsubsection}{\@startsection{subsubsection}{3}{\z@}{-3.25ex\@plus -1ex \@minus -.2ex}{1.5ex \@plus .2ex}{\normalfont\normalsize}}
\makeatother

%%% 04. Tableaux & Débordement Extérieur
%
% MÉMO UTILISATION ORG-MODE (Tableaux) :
%   #+ATTR_LATEX: :environment tabularx :width \largeurimpression :align Z{1.3}Yl :center nil
%   Colonnes : Y = fluide égale | Z{p} = fluide pondérée (ex: Z{1.5}) | l = fixe/naturelle
%
\usepackage{array,booktabs,longtable,tabularx,xltabular,multirow,makecell,diagbox,changepage}

\newcolumntype{Y}{>{\raggedright\arraybackslash}X}
% Z{p} : la pondération modifie \hsize ; \linewidth doit la suivre, sans quoi
% une boîte dimensionnée sur \linewidth dans la cellule (en-tête \eb, liste)
% garde la largeur d'une colonne X moyenne et déborde sur la voisine dès que
% p < 1.
\newcolumntype{Z}[1]{>{\hsize=#1\hsize\linewidth=\hsize\raggedright\arraybackslash}X}

% Contenu des cellules aligné en haut : p{} plutôt que m{}. Avec m{}, une
% cellule courte se centre sur la hauteur de la plus haute de sa ligne, ce qui
% désaligne les tables dont une colonne déborde sur plusieurs lignes.
\renewcommand{\tabularxcolumn}[1]{p{#1}}

% En-tête aligné en bas. Dans une colonne p{}, la ligne de base d'une cellule
% est sa PREMIÈRE ligne : un en-tête court se poserait donc au niveau de la
% première ligne d'un en-tête long. La boîte \parbox[b] prend pour ligne de
% base sa DERNIÈRE ligne, ce qui aligne tous les en-têtes par le bas. À poser
% sur les cellules d'en-tête des colonnes fluides Y et Z{} ; une colonne l ne
% porte qu'une ligne et s'aligne d'elle-même.
\newcommand{\eb}[1]{\parbox[b]{\linewidth}{\raggedright #1}}

\strictpagecheck
% Dans un document recto seul (geometry : twoside=false), la zone d'annotation
% est toujours à droite : aucun décalage n'est alors dû, quelle que soit la
% parité de la page.
\makeatletter
\newcommand{\versexterieur}{\noindent\if@twoside\checkoddpage\ifoddpage\else\hspace*{-\debordementannotation}\fi\fi}
\makeatother
\BeforeBeginEnvironment{tabularx}{\versexterieur}
% longtable et xltabular, qui en dérive, ne composent pas une boîte en ligne :
% chaque rangée est posée entre les ressorts \LTleft et \LTright, à partir de
% la marge gauche du bloc de texte. Le \hspace* de \versexterieur reste donc
% sans effet sur eux, et une table de largeur \largeurimpression composée sur
% une page paire déborde du côté du pli au lieu de la zone d'annotation. On
% règle les deux ressorts à la place, de sorte que la rangée, large de
% \largeurimpression, occupe exactement la largeur du bloc augmentée de la zone
% d'annotation du côté extérieur ; \strictpagecheck rend le test de parité
% sûr dans un flottant, la page étant lue au passage précédent. Le
% \nopagebreak initial retient la table sous la légende que \captionof pose
% juste avant elle lorsque la table n'est pas flottante (:float nil).
%
% Une table longue doit rester hors flottant : dans un flottant, longtable ne
% peut couper la page et sa hauteur échappe au placement. Un flottant reçoit
% donc une table tabularx, une table de plus d'une page une xltabular.
%
% Une xltabular qui franchit une page garde le décalage calculé à son début :
% elle ne peut changer de côté de débordement d'une page à l'autre. Un
% document recto verso dont les tables sont sécables pose donc
% \tableexterieurefalse et les compose dans la largeur du bloc
% (:width \textwidth) ; un document recto seul n'a pas ce problème, sa zone
% d'annotation restant du même côté. Le second \nopagebreak empêche qu'une coupure entre le
% repère de parité et le début de la table fausse le test.
\newif\iftableexterieure
\tableexterieuretrue
\makeatletter
\newcommand{\versexterieurLT}{%
  \ifvmode\nopagebreak\fi
  \iftableexterieure
    \if@twoside\checkoddpage\else\oddpagetrue\fi
    \ifvmode\nopagebreak\fi
    \ifoddpage
      \setlength{\LTleft}{0pt}\setlength{\LTright}{-\debordementannotation}%
    \else
      \setlength{\LTleft}{-\debordementannotation}\setlength{\LTright}{0pt}%
    \fi
  \else
    \setlength{\LTleft}{0pt}\setlength{\LTright}{0pt plus 1fill}%
  \fi}
\makeatother
\BeforeBeginEnvironment{xltabular}{\versexterieurLT}

%%% 05. Flottants, Figures, Tableaux & Légendes
%
% MÉMO UTILISATION ORG-MODE (Figures) :
%   - Figure standard : #+ATTR_LATEX: :width 0.9\linewidth :options alt={Texte alternatif}
%   - Figure large    : #+LATEX: \imagelargetrue
%                       #+ATTR_LATEX: :width \largeurimpression :options alt={Texte alternatif}
%   - Listing         : #+ATTR_LATEX: :caption Légende sous le code
%
% MÉMO UTILISATION ORG-MODE (Description, Note, Source, texte de remplacement) :
%
%   LES QUATRE MOTS-CLÉS PERSONNELS OUVRENT LA PILE. C'est une contrainte d'org,
%   non une préférence : il n'attache à un élément que les mots-clés affiliés qui
%   le précèdent SANS INTERRUPTION, et #+DESC:, #+NOTE:, #+SOURCE: et #+ALT_TEXT:
%   n'en sont pas. Placés après #+CAPTION:, ils coupent la chaîne et la figure
%   perd sa légende, son numéro et son étiquette. Placés en tête, ils ne coupent
%   rien — et le document se compose correctement même si le filtre n'a pas
%   tourné, sans les items mais avec ses légendes.
%
%   #+DESC:     Ce que la figure présente, pour dispenser le texte de le raconter
%   #+NOTE:     Comment lire le graphique lorsque sa forme n'est pas courante
%   #+SOURCE:   D'où vient la donnée ou l'image
%   #+ALT_TEXT: Description destinée à la synthèse vocale
%   #+CAPTION:  Le titre du flottant
%   #+NAME:     fig:etiquette
%   #+ATTR_LATEX: :placement [htbp] :options width=.8\largeurimpression
%   [[./img/figure.pdf]]
%
%   Les quatre premiers sont optionnels et d'ordre libre entre eux.
%   my/org-items-flottants (my-export-config.el) les retire avant l'analyse et
%   les replie là où org les attend : DESC, NOTE et SOURCE dans la légende ;
%   ALT_TEXT dans :options alt={}. Un contrôle refuse toute autre disposition.
%
%   Partage des rôles, à ne pas confondre :
%     #+ALT_TEXT: dit ce que l'image montre                → synthèse vocale
%     #+DESC:     dit ce que l'image présente              → le lecteur
%     #+NOTE:     dit comment lire ce qui est présenté     → le lecteur
%     #+SOURCE:   dit d'où cela vient                      → la provenance
%
\usepackage{graphicx,adjustbox,wrapfig,float}
\usepackage{caption,subcaption}

\captionsetup{format=plain, labelfont=bf, font=small}

\newif\ifimagelarge
\imagelargefalse

%%%% Grille de légende : Titre, Description, Note, Source
%
% Les trois macros ne composent rien : elles déposent leur argument dans un
% registre que le format de légende relit après avoir composé la ligne de titre.
% Le dépôt est global et idempotent, donc insensible au nombre de passes que le
% paquet caption effectue sur la légende.
%
\makeatletter
\newcommand{\itemdescription}{}
\newcommand{\itemnote}{}
\newcommand{\itemsource}{}
\newcommand{\descfig}[1]{\gdef\itemdescription{#1}}
\newcommand{\notefig}[1]{\gdef\itemnote{#1}}
\newcommand{\srcfig}[1]{\gdef\itemsource{#1}}
\newcommand{\razitemsflottant}{%
  \global\let\itemdescription\@empty
  \global\let\itemnote\@empty
  \global\let\itemsource\@empty}
\razitemsflottant

% Colonne d'étiquettes, calibrée sur la plus large une fois les polices chargées
\newlength{\etiquetteflottant}
\newlength{\corpsflottant}
\AtBeginDocument{\settowidth{\etiquetteflottant}{\small\bfseries Description :\hspace{1.4em}}}

% Une ligne de la grille : #1 corps, #2 étiquette, #3 texte
\newcommand{\ligneflottant}[3]{%
  \par\noindent
  \parbox[t]{\etiquetteflottant}{#1\raggedright #2}%
  \parbox[t]{\corpsflottant}{#1\raggedright #3}%
  \par}

% La ligne de titre : #1 largeur totale, #2 étiquette de numéro, #3 séparateur, #4 titre
\newcommand{\grilletitre}[4]{%
  \setlength{\corpsflottant}{\dimexpr#1-\etiquetteflottant\relax}%
  \ligneflottant{\small}{\bfseries#2#3}{#4}}

% Les trois items optionnels : #1 largeur totale
\newcommand{\grilleitems}[1]{%
  \setlength{\corpsflottant}{\dimexpr#1-\etiquetteflottant\relax}%
  \ifx\itemdescription\@empty\else
    \vspace{0.55em}\ligneflottant{\small}{Description :}{\itemdescription}\fi
  \ifx\itemnote\@empty\else
    \vspace{0.55em}\ligneflottant{\footnotesize}{Note :}{\itemnote}\fi
  \ifx\itemsource\@empty\else
    \vspace{0.55em}\ligneflottant{\footnotesize}{Source :}{\itshape\itemsource}\fi}

% Les deux d'un coup, pour un flottant dont la légende est en dessous
\newcommand{\grilleflottant}[4]{\grilletitre{#1}{#2}{#3}{#4}\grilleitems{#1}}

% Vrai si au moins un item a été déposé
\newcommand{\siitemsflottant}[1]{%
  \ifx\itemdescription\@empty\ifx\itemnote\@empty\ifx\itemsource\@empty
  \else#1\fi\else#1\fi\else#1\fi}
\makeatother

%%%% Formats de légende adaptés au débordement extérieur
%
% Tableau : le titre reste au-dessus, où org place la légende ; les trois items
% descendent sous le filet de fin, à la manière d'une note de tableau.
% Figure : la légende étant déjà en dessous, la grille entière y tient.
%
\DeclareCaptionFormat{aplomb}{%
  \versexterieur\parbox{\largeurimpression}{\grilletitre{\largeurimpression}{#1}{#2}{#3}}}

\DeclareCaptionFormat{sousimage}{%
  \ifimagelarge
    \versexterieur\parbox{\largeurimpression}{\grilleflottant{\largeurimpression}{#1}{#2}{#3}}%
  \else
    \grilleflottant{\linewidth}{#1}{#2}{#3}%
  \fi}

\DeclareCaptionFormat{souslecode}{\par\nobreak#1#2#3}

%%%% Application aux types d'entrées
\DeclareCaptionLabelSeparator{grille}{~:}

% Tables : conteneur élargi, titre en aplomb, items sous le tableau
\captionsetup[table]{format=aplomb, singlelinecheck=false, labelsep=grille}
\AtBeginEnvironment{table}{%
  \setlength{\linewidth}{\largeurimpression}%
  \setlength{\hsize}{\largeurimpression}%
  \setlength{\columnwidth}{\largeurimpression}}
\AtEndEnvironment{table}{%
  \siitemsflottant{%
    \par\addvspace{0.7em}%
    \versexterieur\parbox{\largeurimpression}{\grilleitems{\largeurimpression}}\par}%
  \razitemsflottant}

% Figures : format dynamique avec réinitialisation automatique des registres
\captionsetup[figure]{format=sousimage, singlelinecheck=false, labelsep=grille}
\AtEndEnvironment{figure}{\global\imagelargefalse\razitemsflottant}

% Listings : légende en dessous, typographie dédiée
\captionsetup[listing]{format=souslecode, labelfont=bf, textfont=it}

%%% 06. Blocs de code & Listes standard
\usepackage[normalem]{ulem}
\usepackage{enumitem}\setlist{nosep}\setlist[itemize]{leftmargin=*}
\usepackage{fvextra} % Requis avant csquotes
\usepackage{csquotes}
\usepackage{newfloat}
\DeclareFloatingEnvironment[fileext=lst,listname={},name=Listing]{listing}
\fvset{breaklines=true,breakanywhere=true}

%%% 07. Factorisation des listes & Messages de repli
%
% MÉMO UTILISATION :
%   \listoffigures, \listoftables, \listoflistings affichent un message bilingue si vides.
%   Pour brancher une nouvelle liste (ex: glossaire) :
%     \renewcommand{\printglossary}{\ListeOuRepli{glo}{gls}{\videglossaire}}
%
\newcommand{\FrEn}[2]{\ifdefstring{\languagename}{french}{#1}{#2}}
\newcommand{\MessageVide}[2]{%
  \par\smallskip\noindent\emph{\FrEn{#1 dans ce document.}{#2 in this document.}}\par\smallskip}
\newcommand{\DeclareRepliVide}[3]{\csdef{vide#1}{\MessageVide{#2}{#3}}}

\DeclareRepliVide{figures}  {Aucune figure}              {No figure}
\DeclareRepliVide{tableaux} {Aucun tableau}              {No table}
\DeclareRepliVide{listings} {Aucun listing}              {No listing}
\DeclareRepliVide{index}    {Aucune entrée d'index}      {No index entry}
\DeclareRepliVide{acronymes}{Aucun acronyme}             {No acronym}
\DeclareRepliVide{glossaire}{Aucune entrée de glossaire} {No glossary entry}

\makeatletter
\newcommand{\ListeOuRepli}[3]{\ifnum\value{#1}>\z@ \@starttoc{#2}\else #3\fi}
\renewcommand{\listoffigures}{\ListeOuRepli{figure}{lof}{\videfigures}}
\renewcommand{\listoftables}{\ListeOuRepli{table}{lot}{\videtableaux}}
\renewcommand{\listoflistings}{\ListeOuRepli{listing}{lst}{\videlistings}}
\makeatother

%%% 08. Notes marginales, Annexes & Bibliographie
%
% MÉMO UTILISATION ORG-MODE :
%   - Remarque marginale : \RMQ{Texte de la note en marge} (balisée sémantiquement en <Aside>)
%   - Bibliographie      : #+print_bibliography: :heading subbibliography
%
\usepackage[toc,page]{appendix}
\usepackage{marginnote}
\usepackage{perpage}\MakePerPage{footnote}
\renewcommand{\thefootnote}{\alph{footnote}}
\renewcommand*{\marginfont}{\footnotesize\raggedright}

% Remarque numérotée conforme PDF/UA-2 (Structure Aside)
\newcounter{rmq}
\renewcommand{\thermq}{\arabic{rmq}}
\newcommand{\RMQ}[1]{%
  \refstepcounter{rmq}%
  \tagstructbegin{tag=Aside}%
    \marginnote{\textbf{RMQ~\thermq.}~#1}%
  \tagstructend%
  \textsuperscript{\textsf{R\thermq}}}

\usepackage[
  bibstyle=iso-numeric,
  citestyle=numeric-comp,
  doi=true,
  isbn=false,
  mincrossrefs=2,
  sorting=none,
  sortcites=true,
  autolang=other,
  backend=biber]{biblatex}
%
% refsection=section retiré le 8 septembre. L'option demandait à biblatex de
% rustiner \@startsection, \@sect et \@ssect au \begin{document} pour ouvrir
% une refsection à chaque titre. Les trois rustines échouaient, titletoc (§ 10)
% ayant déjà redéfini ces trois macros — trois erreurs à chaque compilation :
%
%     Package biblatex Error: Patching \@startsection failed.
%
% L'option était de surcroît sans emploi : le crochet du § 10 ouvre déjà la
% refsection à la main, et c'est lui qui donne aux sous-bibliographies leur
% segment. Retirer l'option ne change donc rien au rendu, et rend le journal
% de compilation lisible.
\renewcommand*{\bibfont}{\scriptsize}

%%% 09. Environnements Mathématiques & Théorèmes
%
% MÉMO UTILISATION ORG-MODE :
%   #+BEGIN_EXPORT latex
%   \begin{theorem}[Titre optionnel]
%     \begin{statement} Énoncé \end{statement}
%     \begin{proofsketch} Idée de la démonstration \end{proofsketch}
%   \end{theorem}
%   #+END_EXPORT
%
\AtBeginDocument{%
  \let\llbracket\lBrack \let\rrbracket\rBrack
  \let\Box\mdlgwhtsquare \let\square\mdlgwhtsquare
  \let\Diamond\mdlgwhtdiamond \let\bigcirc\mdlgwhtcircle
  \let\leadsto\rightsquigarrow
}

\newcounter{theoreme}
\newcommand{\entetetheoreme}[3]{\par\addvspace{0.6em}\noindent{#1#2\ifblank{#3}{}{ : #3}}\par\nobreak\vspace{0.15em}\nobreak}
\newenvironment{theorem}[1][]{\refstepcounter{theoreme}\entetetheoreme{\bfseries}{Théorème \thetheoreme}{#1}}{\par\addvspace{0.6em}}
\newenvironment{statement}[1][]{\entetetheoreme{\normalfont}{Déclaration \thetheoreme}{#1}}{\par}
\newenvironment{proofsketch}{\entetetheoreme{\itshape}{Esquisse de preuve}{}\itshape\ignorespaces}{\unskip\nobreak\hfill$\mdlgwhtsquare$\par}

%%% 10. Pagination & Tables des matières (ToC)
%
% MÉMO UTILISATION ORG-MODE :
%   - ToC générale 1er niveau : #+LATEX: \clearpage
%                               #+TOC: headlines 1
%                               #+LATEX: \clearpage
%   - ToC locale de section   : #+TOC: headlines 2 local
%
\usepackage{lastpage}
\fancyfoot[C]{\thepage\ / \pageref{LastPage}}

% Rupture de page et nouvelle section biblio via hook officiel (compatible balisage PDF/UA)
\AddToHook{cmd/section/before}{\cleardoublepage\newrefsection}

\usepackage{titletoc}
\titlecontents{subsection}[1.5em]{\small}{\contentslabel{2.3em}}{\hspace*{-2.3em}}{\titlerule*[1pc]{.}\contentspage}
\titlecontents{subsubsection}[3.8em]{\small}{\contentslabel{3.2em}}{\hspace*{-3.2em}}{\titlerule*[1pc]{.}\contentspage}

%%% 11. ORCID & Hyperliens (Toujours en dernier)
%
% MÉMO UTILISATION ORG-MODE :
%   - Lien ORCID : \orcidlink{0000-0000-0000-0000}
%
\newcommand{\orcidicon}{~/.emacs.d/assets/ORCID-iD-icon-BW-16x16.png}
\newcommand{\orcidlink}[1]{\href{https://orcid.org/#1}{\raisebox{-0.1em}{\includegraphics[height=\fontcharht\font`A]{\orcidicon}}}}

\usepackage{hyperref}
\pdfstringdefDisableCommands{%
  \let\par\empty \let\\\empty
  \def\footnote#1{} \def\marginpar#1{} \def\marginnote#1{}%
  \def\href#1#2{#2} \def\textcolor#1#2{#2} \def\uline#1{#1} \def\sout#1{#1} \def\enquote#1{#1}%
  \def\orcidlink#1{ORCID: #1}%
}
\hypersetup{
  colorlinks=true,
  linkcolor=customgray,
  citecolor=customgray,
  urlcolor=customgray,
  pdfborder={0 0 0},
  unicode=true,
  breaklinks=true,
  bookmarksnumbered=true,
  pdflang=fr,
  pdfdisplaydoctitle=true}
\usepackage{bookmark}
